\documentclass[11pt,letterpaper]{article}
\usepackage[T1]{fontenc}
\usepackage[utf8]{inputenc}
\usepackage{lmodern,microtype}
\usepackage[margin=.94in,headheight=15pt]{geometry}
\usepackage{amsmath,amssymb,amsthm,mathtools}
\usepackage{booktabs,longtable,array,enumitem,xcolor,fancyhdr,listings,needspace}
\usepackage[hyphens]{url}
\usepackage{hyperref}
\hypersetup{colorlinks=true,linkcolor=blue!45!black,urlcolor=blue!45!black,
 citecolor=blue!45!black,pdftitle={Completing the Successor-Seeded Scalar Basis for iPRA},
 pdfauthor={Research prepared for Vladimir Reshetnikov}}
\setlength{\emergencystretch}{3em}
\setlength{\parskip}{.22em}
\setlist{itemsep=.2em,topsep=.4em,leftmargin=1.7em}
\newtheorem{theorem}{Theorem}[section]
\newtheorem{lemma}[theorem]{Lemma}
\newtheorem{proposition}[theorem]{Proposition}
\newtheorem{corollary}[theorem]{Corollary}
\theoremstyle{definition}\newtheorem{definition}[theorem]{Definition}
\theoremstyle{remark}\newtheorem{remark}[theorem]{Remark}
\newcommand{\ipra}{\mathsf{iPRA}}
\newcommand{\D}{\mathsf D}
\newcommand{\Q}{\mathsf Q}
\newcommand{\FV}{\operatorname{FV}}
\newcommand{\file}[1]{{\small\nolinkurl{#1}}}
\lstset{basicstyle=\ttfamily\footnotesize,columns=fullflexible,keepspaces=true,
 breaklines=true,showstringspaces=false,frame=single,rulecolor=\color{black!20},
 xleftmargin=.3em,xrightmargin=.3em}
\pagestyle{fancy}\fancyhf{}
\fancyhead[L]{\small Successor-seeded scalar arithmetic}
\fancyhead[R]{\small Construction, equivalence, and limits}
\fancyfoot[C]{\thepage}
\title{\textbf{Completing the Successor-Seeded\\Scalar Basis for iPRA}\\[.5em]
\large A predecessor construction resolving an open candidate in the supplied drafts}
\author{Research prepared for Vladimir Reshetnikov}
\date{November 10, 2026}
\begin{document}
\maketitle
\begin{abstract}
The shared-continuation draft proposes a numerical language consisting of
zero and a scalar iterator that starts at the successor of its initializer.
It explicitly leaves open the construction of a predecessor in this
language. We complete that candidate. A shifted factor map yields an
internally correct comparator; a saturated iteration yields a capped
predecessor; a fixed three-input correction yields a genuine predecessor.
Every arithmetic induction in the new bootstrap has a single literal
equation as its invariant. No predecessor, successor cancellation, ordinary
initialization, pairing, or general selector is used to manufacture the
predecessor.

The resulting basis has two numerical constructors, two computation
schemas, and one changing update variable. With the explicit native
logical rules given below it is bi-interpretable and equiconsistent with
full first-order intuitionistic primitive recursive arithmetic. Both
interpretations use the entire numerical domain and actual equality, and
both numerical composite isomorphisms are the identity.

The mathematical completion and the implementation comparison are separate
claims. The supplied complete checker improves the inspected byte frontier,
but does not establish a simultaneous improvement in every source-size
coordinate, absolute minimality, or global priority. The equivalence proof
is a written proof with independent assistant review and executable
diagnostics; it is not a Lean or Rocq formalization.
\end{abstract}
\clearpage
\begingroup
\small\setlength{\parskip}{0pt}
\setcounter{tocdepth}{2}\tableofcontents
\endgroup
\clearpage

\section{What is new, and what has actually been minimized}
\label{sec:result}

\subsection{The resolved obligation}
The comparison is pinned to ProveIt commit
\begin{center}\file{2fdcf5bf988bead821c416649f713f98b7ed8776}.\end{center}
The current directory contains nine direct ZIP archives. In
\file{ipra_shared_continuation_2026-11-10.zip}, the report's section
``A genuinely smaller numerical grammar: the precise unresolved step''
states exactly the successor-seeded scalar candidate studied here and
proves a conditional initialized-iterator compiler. Its unresolved
obligation is a finite term $P$ and an internal proof of $P(Sa)=a$
\cite{shared}. The present work supplies that term and proof, including
$P(0)=0$. It therefore completes an existing proposed basis rather than
claiming to have invented its syntax.

The new steps are the shifted affine alignment of the factor-map
construction, its whole-domain comparator proof, and the finite correction
of a capped predecessor. The earlier ordinary scalar bootstrap already
uses finite cycles, a factor map, and saturation; that strategy is credited
to the supplied development\cite{generated}. Related restricted-iteration
and seed-shifting methods occur in Severin\cite{severin}. This limited
literature check establishes context, not a global priority claim.

\begin{center}
\begin{tabular}{lccc}
\toprule
Numerical basis & Constructors & Computation schemas & Update variables\\
\midrule
Initialized scalar $0,S,I$ & 3 & 2 & 1\\
Counter-aware successor $0,R$ & 2 & 2 & 2\\
Completed scalar $0,J$ & 2 & 2 & 1\\
\bottomrule
\end{tabular}
\end{center}
The counts omit variables and count a binder-forming operation as one
constructor. The $R$ comparison is the existing successor-recursion
proposal\cite{successor}. Two numerical constructors alone are not a
proof-checker lower bound: argument structure, logic, validation, and
encoding also cost code.

\subsection{Complete-checker comparison}
\begin{center}
\begin{tabular}{lrrr}
\toprule
Complete executable & Tokens & AST nodes & Normalized bytes\\
\midrule
Prior generated-assumption checker & 281 & 318 & 685\\
Prior AST-oriented checker & 283 & 315 & 710\\
New $J$: direct adaptation & 283 & 320 & 687\\
New $J$: byte profile & 284 & 322 & \textbf{681}\\
New $J$: AST profile & 282 & \textbf{313} & 707\\
New $J$: total-input AST profile & 287 & \textbf{314} & 711\\
\bottomrule
\end{tabular}
\end{center}
The new separate-profile envelope is $282/313/681$, compared with the
prior $281/315/685$. It improves AST size by two nodes and normalized
length by four bytes. It does not improve the token minimum. The AST
profile's additional raw aliases and its reliance on CPython's finite
recursion limit for malformed nonempty strings are documented in
Section~\ref{sec:implementation}.
The additional total-input profile improves AST size by one node while
retaining termination on every finite malformed input in the ideal
unbounded-resource model. Its separate-profile envelope with the direct
and byte versions is $283/314/681$.

The audit recursively includes embedded source archives and complete
programs stored in reference directories. It found 415 Python source
occurrences, 164 distinct source hashes, and no parse errors. Fourteen
distinct smaller-coordinate utilities were inspected and excluded because
they are builders, measurers, integrity verifiers, or test runners calling
another checker. The complete measured frontier is consequently
$281$ normalized tokens, $315$ AST nodes, and $685$ normalized bytes;
different existing programs attain these coordinates.

The metric removes comments and actual docstrings, applies
\texttt{ast.unparse}, and measures the entire module, including JSON input,
scope checks, inference checks, rejection, and verdict output. Python and
its standard JSON parser are common runtime infrastructure. The supplied
normalizer and frozen controls reproduce the comparison under CPython
3.12.14. An observed minimum among implementations is an attained upper
bound on the possible checker cost, not a proved mathematical lower bound.

The byte saving is small and partly a Python source representation change.
The substantive mathematical result is the resolved predecessor
obligation. We do not attribute a representation-only saving to the
arithmetic theorem, nor claim that a single supplied program beats all
three current coordinate minima.

\section{The entire primitive theory}
\label{sec:primitive}

Here $\ipra$ is the full first-order intuitionistic theory of fixed finite
primitive-recursive descriptions, their equations, the successor boundary
laws, and quantifier-free induction. The atomic-induction presentation is
standard\cite{salehi}; the explicit equality-reflection argument in
Appendix~\ref{sec:eic-completion} verifies the relevant equivalence here.
It is important to retain the full first-order logical interface.
Equivalence only with the quantifier-free equational fragment would not
satisfy the requested bi-interpretability condition.

\subsection{Numerical language}
The finite staged terms are
\[
 t::=x\mid0\mid J(n,a;u.h),\qquad Sa:=J(0,a;u.0),\qquad1:=S0.
\]
Only $u$ in the update $h$ is bound. Its remaining parameters are fixed
throughout that iteration. The numerical axioms, available in any retained
context, are exactly
\[
 J(0,a;u.h)=Sa,\qquad
 J(Sn,a;u.h)=h[J(n,a;u.h)/u].\tag{J0, JS}
\]
Every substitution is capture-free. Equivalently introduce one ordinary
function symbol $j_h(n,a,\vec p)$ for each fixed earlier-stage description
$h(u,\vec p)$. In particular the zero-body description defines successor.
There is no object-level quantification over descriptions, no higher-type
function sort, and no total universal evaluator for all PR descriptions.

\subsection{Native formulas and every logical rule}
Formulas and reusable contexts have grammar
\[
 A::=(s=t)\mid\Q_y A\mid\D_{x,y}A,\qquad
 \Gamma\text{ a finite reusable list of such formulas}.
\]
The variables $x,y$ are distinct. Both are bound by $\D$; only $y$ is
bound by $\Q$. Their intended ordinary meanings are
\[
 (\Q_yB)^\circ=B^\circ(0)\to\forall yB^\circ(y),
 \qquad
 (\D_{x,y}B)^\circ=
 \exists x\bigl(B^\circ(x,0)\to\forall yB^\circ(x,y)\bigr).
\]
Besides the two computation schemas, the complete primitive inference
inventory is
\[
 \frac{}{\Gamma\vdash A}\ (A\in\Gamma),\qquad
 \frac{\Gamma\vdash s=t\quad\Gamma\vdash F(s)}
      {\Gamma\vdash F(t)},\qquad
 \frac{\Gamma\vdash\Q_yB\quad\Gamma\vdash B(0)}
      {\Gamma\vdash B(t)}.
 \tag{Hyp, Tr, QE}
\]
The transport motive $F$ is any fixed native formula template. The paired
elimination and its two introduction choices are
\[
 \frac{\Gamma\vdash\D_{x,y}B(x,y)\quad
       \Gamma,\Q_yB(a,y)\vdash F}{\Gamma\vdash F},
 \tag{DE}
\]
\[
 \frac{\Gamma,B(t,0)\vdash B(t,Sz)}{\Gamma\vdash\D_{x,y}B(x,y)},
 \qquad
 \frac{\Gamma,E(t,z)\vdash E(t,Sz)}{\Gamma\vdash\D_{x,y}E(x,y)}.
 \tag{DI-case, DI-progress}
\]
The matrix $E$ in DI-progress must be a literal numerical equation.
An equation under either binder does not qualify. In DI, $z$ is absent
from the retained context, witness, and passive matrix parameters.
In DE, $a$ is absent from the retained context, conclusion, and passive
matrix parameters. The displayed local assumption is discharged.
Bound variables and eigenparameters are renamed before substitution.
These rules do not contain a classical rule or induction on a compound
formula. Native falsity is the formula $0=1$.

Weakening, exchange, contraction, numerical substitution, renaming, and
ordinary assumption cut are admissible by induction on derivations.
Rename local eigenparameters first. Numerical substitution preserves the
root constructor of an equation, so it preserves the progress restriction.
For cut, replace each use of the discharged Hyp by its given proof and
recurse through the remaining rule instances. These are metatheorems,
not uncharged primitive instructions in the checker.

\subsection{The equivalent computation-supply presentation}
The implementation uses the inherited generated-assumption presentation
\cite{generated}. It replaces a numerical axiom leaf by
\[
 \frac{\Gamma,H(w)\vdash A}{\Gamma\vdash A},\qquad
 w\notin\FV(\Gamma,A),\quad \Gamma\vdash H(0).
 \tag{Known-instance supply}
\]
The checker generates only the particular equations whose zero instances
are justified by J0 or JS and the raw opening lemma; it does not decide
arbitrary instances of the displayed side condition. The rule is
admissible: substitute zero for $w$ in the premise and cut the proof of
$H(0)$. Conversely, a computation theorem $E$ is obtained by supplying
the constant equation $H(w)=E$ and immediately using Hyp. Thus the
computation-leaf and computation-supply versions prove the same sequents.
This proof architecture is inherited, not a new contribution here.

% Native bridge audit and the pre-selector finite-case repair.
% Requires amsmath, amssymb, amsthm and theorem/lemma environments.

\section{Native bootstrap: cases before a full selector}
\label{nb:section}

This section checks a dependency that matters for the new scalar basis.
Before predecessor is available, the partial selector
\[
 B(n,a,b)=J(n,a;u.b),\qquad B(0,a,b)=Sa,\quad B(Sn,a,b)=b
\]
is not the unrestricted selector $C(0,a,b)=a$, $C(Sn,a,b)=b$.
Accordingly, one must not silently invoke the later ordinary-connective
macros during the arithmetic bootstrap.  In particular, the theorem rule
for numerical cases is not by itself a rule for eliminating equations
$t=0$ and $t=Sz$ from a retained context.  The finite-case lemma below
supplies exactly the stronger form used by the bootstrap.

We use the native Eq/$Q$/$D$ presentation with hypothesis use, full
formula transport, $Q$ elimination, paired elimination, and the two
paired introductions.  In notation,
\[
 Q_yF=\text{the clause taking a proof of }F(0)\text{ to }F(t),
\]
\[
 \frac{\Gamma,F(t,0)\vdash F(t,Sz)}
      {\Gamma\vdash D_{x,y}F(x,y)},\qquad
 \frac{\Gamma,E(t,z)\vdash E(t,Sz)}
      {\Gamma\vdash D_{x,y}E(x,y)}.
\]
The second matrix is a literal equation.  The usual freshness conditions
are retained, and paired elimination opens $Q_yF(a,y)$ at a fresh $a$.
Computation equations may be read as leaves, or obtained from their
constant generated-assumption instances and hypothesis use.

\subsection{Resources available without a numerical selector}

Choose $v$ fresh for $t$ and put $r=J(S0,0;v.t)$.  Computation gives
$r=t$.  Transport this equation along itself with motive $w=t$ to obtain
$t=t$.  The familiar equality motives then give symmetry, transitivity
and numerical congruence; this uses neither predecessor nor successor
injectivity.

For any native formula template $F$, case introduction with an unused
witness coordinate gives $D_{x,y}F(y)$ from
$\Gamma\vdash F(Sz)$, ignoring the retained $F(0)$.  Paired elimination
opens $Q_yF(y)$; a supplied proof of $F(0)$ followed by $Q$ elimination
at $t$ gives
\[
 \frac{\Gamma\vdash F(0)\qquad\Gamma\vdash F(Sz)}
      {\Gamma\vdash F(t)}. \tag{Native NC}\label{nb:nc}
\]
The identical construction with progress introduction, using its local
$E(z)$, gives contextual literal-equation induction from a base and step.
These arguments use no arithmetic selector or ordinary logical connective.

For reference, $Q_yF$ and $D_{x,y}F$, with $x$ unused, are interderivable
already: use case introduction and $Q$ elimination in the forward
direction, and paired elimination followed by hypothesis use in the
reverse direction.  Thus the additional $Q$ constructor presents the
same internal clauses as the Eq/$D$ formulation.

Even equation explosion is available at this stage.  Under $0=1$,
congruence gives $B(0,0,t)=B(1,0,t)$; computation gives $1=t$.
Consequently every equation follows.  External induction on fixed native
formula syntax, using case introduction (and the unused-witness
representation for $Q$), gives explosion for every native formula.
This is an external proof transformation, not compound-formula induction
inside the arithmetic theory.

\subsection{Functional finite ranges do not assume branch equations}

Say that a particular term $b$ has the \emph{functional range rule}
$\operatorname{FR}_R(b)$, for a finite list of numerals $R$, when the
following metarule is derived for every native formula template $K$:
\[
 \bigl(\Gamma\vdash K(r)\text{ for each }r\in R\bigr)
 \quad\Longrightarrow\quad\Gamma\vdash K(b).
 \tag{FR}\label{nb:fr}
\]
There are no added equations $b=r$ in these premises.  The retained
context and every fixed parameter of $K$ remain unchanged.  In
particular, they may mention the original term $b$.

Use the fixed terms $Z,N,j,M$ defined in
Section~\ref{js:cycles}, with $Z(t)=B(t,0,0)$. The present proof
checks their logical prerequisites; their finite computation identities
are verified directly from those definitions.

For $Z(t)$, apply Native NC to the fresh template $w\mapsto K(Z(w))$.
The zero branch is $K(1)$ and every successor branch is $K(0)$, by
computation and transport.  Hence $\operatorname{FR}_{\{0,1\}}(Z(t))$
holds, and therefore so does the same rule for $N(t)$.
For $B(t,1,0)$ the identical construction has endpoints two and zero.
For
\[
 j(t)=B\bigl(t,0,B(Z(N(t)),1,0)\bigr)
\]
use NC on the outer control: the zero branch is one, and the successor
branch has the preceding two-value rule.  This proves
$\operatorname{FR}_{\{0,1,2\}}(j(t))$, and hence the same rule for
$M(t)$, since $M$ is an outer application of $j$.

The inner parity iterate $U(t)=J(t,0;u.Z(u))$ has a two-value range
rule as well.  At counter zero it is one; at a successor it is a value
of $Z$.  Apply NC on the counter and the already proved range rule of
$Z$.  No induction with a disjunctive invariant is involved.

These constructions justify the cycle equations before contextual
finite-case elimination is used.  For example, $Z(Z(U(t)))=U(t)$ follows
by applying its two-value rule to the equation template $Z(Z(w))=w$.
The equations
\[
 M(0)=0,\quad M(1)=1,\quad M(2)=2,\qquad M(Sw)=j(M(w))
\]
and $j(0)=1$, $j(1)=2$, $j(2)=0$ then follow from the fixed definitions.
Their functional range proofs are independent of the ramp invariant.

\begin{lemma}[Successor cancellation on the three finite values]
\label{nb:finite-cancellation}
Suppose $b$ and $c$ have functional range rules whose values are contained
in $\{0,1,2\}$.  In any retained context,
\[
                         Sb=Sc\quad\Longrightarrow\quad b=c
\]
is derivable, before global successor injectivity is known.
\end{lemma}
\begin{proof}
Apply the two functional range rules to the equation templates
$M(w)=w$ and $j(j(j(w)))=w$.  Their required numeral instances are fixed
computations.  Thus $M(b)=b$, $M(c)=c$, $j^3(b)=b$ and $j^3(c)=c$
are already theorems in the retained context.
Apply $M$ to $Sb=Sc$ and use $M(Sw)=j(M(w))$ to get $j(b)=j(c)$.
Apply $j$ twice and use the two cycle identities to obtain $b=c$.
Every step is ordinary numerical congruence or transport.  In particular,
no cancellation of an arbitrary successor was inserted.
\end{proof}

\begin{lemma}[Contextual finite-case elimination for equations]
\label{nb:guarded-cases}
Let $b$ have $\operatorname{FR}_R(b)$, where $R$ is contained in
$\{0,1,2\}$, and let $F$ be the literal equation $u=v$.
If, for each $r\in R$, there is a derivation
\[
                         \Gamma,b=r\vdash u=v,
\]
then $\Gamma\vdash u=v$ is derivable.  The context may contain arbitrary
native formulas, including an induction hypothesis that depends on $b$.
\end{lemma}
\begin{proof}
Define a native formula template, with both bound variables fresh,
\[
 K(x):=D_{a,y}\bigl(B(y,x,u)=B(y,b,v)\bigr),
\]
where $a$ is unused.  Fix $r\in R$.  To introduce $K(r)$ by the case
choice with witness zero, its local zero assumption reduces by
computation to $Sr=Sb$.  Lemma~\ref{nb:finite-cancellation}, with the
constant $r$ and the known finite-range term $b$, yields $b=r$.
Cut in the given branch derivation.  Its conclusion $u=v$ is exactly
the successor matrix, up to the computations of $B$.  Thus
$\Gamma\vdash K(r)$.

Apply the functional range rule of $b$ to these closed branch proofs.
It gives $\Gamma\vdash K(b)$, whose matrix is
$B(y,b,u)=B(y,b,v)$.  Paired elimination opens its $Q$ clause.
At zero its two sides are $Sb$, so reflexivity and computation supply
the base.  Specializing the clause at one gives $u=v$.  Close paired
elimination.  All temporary coordinates are fresh and all temporary
assumptions have been discharged.
\end{proof}

The derivation is a finite guard construction, not a primitive implication
connective.  It relies only on the partial selector, computations,
transport, hypotheses, case introduction, paired elimination and $Q$
elimination.  Neither predecessor nor the unrestricted selector occurs.

\subsection{Audit of the bootstrap and completion of the bridge}

The dependency order is now explicit.  First establish functional range
rules and the two finite cycles using only Native NC and equations.
Then establish finite successor cancellation and
Lemma~\ref{nb:guarded-cases}.  The finite splits in the ramp induction,
the affine scaling induction and the factor-map calculation are instances
of that lemma, with $b=N(n)$ or another displayed finite-range term and
with a literal-equation conclusion.  Retained induction hypotheses are
therefore legitimate.  The saturation identity $K(c,Su)=S K(c,u)$ is
proved directly by NC with no branch assumptions; it does not require
an unrestricted contextual split on arbitrary $c$.  No step in this order
uses the later ordinary logical macros.

After the corrected predecessor $P$ is constructed, the initializer
compiler gives the ordinary scalar $I$.  Consequently
\[
 C(n,a,b):=I(n,a;u.b)
\]
has the unrestricted selector endpoints.  The inherited gate, padding,
formula-selector and ordinary-connective constructions may now be applied
in their stated order.  Binder congruence itself uses only case
introduction, paired elimination, transport and cut; it is already valid
before this final stage.

For clarity, the ordinary interpretations of the native binders are
\[
 (Q_yB)^\circ=B^\circ(0)\to\forall y\,B^\circ(y),\qquad
 (D_{x,y}B)^\circ=
 \exists x\bigl(B^\circ(x,0)\to\forall y\,B^\circ(x,y)\bigr).
\]
Once the logical macros are available, the native return translations
are proved by the same explicit constructions:
\[
 Q_yB\dashv\vdash B(0)\Rightarrow_*\forall_*y\,B(y),
\]
\[
 D_{x,y}B\dashv\vdash
 \exists_*x\bigl(B(x,0)\Rightarrow_*\forall_*y\,B(x,y)\bigr).
\]
Forwards, an opened $Q$ clause under its base supplies the arbitrary
instance needed by universal introduction.  Backwards, a displayed
implication and universal supply the successor matrix for case
introduction; an unused witness gives the $Q$ case.  The existential
witness is retained in the $D$ case.  External structural induction on
finite formulas, together with the endpoint proofs for the ordinary
macros, yields both formula composites.  Falsity uses the explicit
ordinary separation bridge $(0=1)\to\bot$.

Finally, the inherited scalar completion and fixed-description PR
compiler are used only after the logical interface and initialized $I$
have been recovered.  Recurrence uniqueness gives the numerical
composites on the whole original domain.  The domain is unchanged,
interpreted equality is actual equality, and the numerical isomorphism
graph is $x=y$.  Thus the finite-case repair closes the pre-selector
logical dependency; it does not replace the separately supplied complete
arithmetic and logical interpretations by standard-model calculations.

% Successor-seeded scalar iteration: a self-contained arithmetic bootstrap.
% Requires amsmath, amssymb, amsthm and theorem/lemma/proposition environments.
% No document wrapper. All induction invariants below are literal equations.

\section{A successor-seeded scalar basis with two computation schemas}
\label{js:section}

This section completes the explicitly unresolved one-variable
successor-seeded candidate in the supplied shared-continuation
report~\cite{shared}, in its section ``A genuinely smaller numerical
grammar: the precise unresolved step.''  The candidate basis and its
conditional initializer compiler are inherited from that report.
Unlike the counter-aware successor recursor, its update has just one
changing input.  The initializer remains a separate argument, but the
iterator starts at its successor.  The missing result is a numerical
predecessor term with a whole-domain proof, obtained here by adapting
the supplied scalar $2/3$ factor-map construction to shifted arithmetic
and correcting a saturated predecessor at three small arguments.

\subsection{Syntax, computations, and the precise resources used}

The terms are finite staged expressions
\[
 t ::= x\mid 0\mid J(n,a;u.h).
\]
Only $h$ binds $u$; its other free variables are passive parameters.
Define the abbreviation
\[
 Sa:=J(0,a;u.0),\qquad 1:=S0,
\]
and use standard numerals as finite strings of this abbreviation.
There are exactly two numerical computation schemas:
\begin{align}
 J(0,a;u.h)&=Sa,\tag{J0}\label{js:j0}\\
 J(Sn,a;u.h)&=h[J(n,a;u.h)/u].\tag{JS}\label{js:js}
\end{align}
The first schema, with the abbreviation expanded, says that
$J(0,a;u.h)=J(0,a;u.0)$.  It does not presuppose a separate successor
constructor.  No function variables or uniform evaluator are present.

The bootstrap uses equality replacement, numerical cases
\[
 \frac{\Gamma\vdash F(0)\qquad\Gamma\vdash F(Sz)}
      {\Gamma\vdash F(t)},
 \tag{NC}\label{js:nc}
\]
and contextual induction only on a literal equation $E$:
\[
 \frac{\Gamma\vdash E(0)\qquad\Gamma,E(z)\vdash E(Sz)}
      {\Gamma\vdash E(t)}.
 \tag{EI}\label{js:ei}
\]
The ordinary eigenvariable conditions apply.  In an ordinary
intuitionistic first-order presentation, NC follows from coverage
$\forall x\,(x=0\vee\exists y\,x=Sy)$.  In the native paired-binder
presentation it is the already available case rule.  No use of
successor injectivity, predecessor, pairing, an equality characteristic,
or induction on a compound formula is made before predecessor is
constructed below.

The finite range splits used below require the distinction proved in
Section~\ref{nb:section}. Functional range elimination follows from NC;
contextual finite cases with retained induction hypotheses use the finite
cancellation and guard construction of Lemma~\ref{nb:guarded-cases}.
The cycle identities are established first by functional range elimination
alone. Thus a displayed finite range is not an uncharged disjunction or
an assumption of the ordinary-connective macros.

We shall use the following admissible consequence of EI.  If two fixed
updates are pointwise equal, and the displayed initial values agree,
their iterates agree at every counter.  Induct on the literal equation
between the two iterates: J0 supplies its base and JS supplies its step.
This is a derived uniqueness argument, not a primitive extensionality
rule for update bodies.  All passive arguments are kept fixed while
performing the induction.

\subsection{A partial selector and shifted addition}

Put
\[
 B(n,a,b):=J(n,a;u.b),\qquad Z(n):=B(n,0,0),
 \qquad a\oplus b:=J(b,a;u.Su).
\]
Their computation endpoints are
\[
 B(0,a,b)=Sa,\quad B(Sn,a,b)=b,\qquad
 Z(0)=1,\quad Z(Sn)=0,
\]
\[
 a\oplus0=Sa,\qquad a\oplus Sb=S(a\oplus b).
\]
The operation $\oplus$ has standard meaning $a+b+1$.  Ordinary
unshifted addition has not been defined and is not used in the proof.
Literal-equation induction proves
\[
 0\oplus a=Sa,\quad Sa\oplus b=S(a\oplus b),\quad
 a\oplus b=b\oplus a,\quad
 (a\oplus b)\oplus c=a\oplus(b\oplus c).
 \tag{Shifted algebra}\label{js:oplus}
\]
For the first two equations induct on the second argument.  For
commutativity induct on $b$, using the first two equations.  For
associativity induct on $c$; at zero both sides are
$S(a\oplus b)$, and the successor step applies $S$ to the previous
equality.  These proofs do not cancel a successor.

In particular, the following are genuine available terms:
\[
 D(a):=S^2(a\oplus a),\qquad
 T(a):=S^4((a\oplus a)\oplus a),\qquad
 R(a):=S^3(a\oplus a).
 \tag{Affine terms}\label{js:affine}
\]
Their standard meanings are respectively $2a+3$, $3a+6$, and
$2a+4$.  The proof uses their displayed term definitions.  From
\eqref{js:oplus} one obtains
\[
 D(0)=3,\quad D(Sa)=S^2D(a),\qquad
 T(0)=6,\quad T(Sa)=S^3T(a).
 \tag{Affine shifts}\label{js:affine-shifts}
\]

\subsection{Two finite cycles}\label{js:cycles}

Define
\[
 N(n):=Z\bigl(J(n,0;u.Z(u))\bigr).
\]
The inner iterator is always zero or one: at zero it is one, and at a
successor it is a value of $Z$.  The endpoints of $Z$ and finite cases
therefore give
\[
 N(0)=0,\quad N(Sn)=Z(N(n)),\quad N(S^2n)=N(n),
 \qquad N(n)\in\{0,1\}.
 \tag{Parity cycle}\label{js:parity}
\]
This asserts a proved two-state recurrence, not a pre-existing parity
operation.  For example, if $K(n)=J(n,0;u.Z(u))$, then
$N(Sn)=Z(Z(K(n)))=K(n)=Z(N(n))$ by its two cases.

Next put
\[
 j(u):=B\bigl(u,0,B(Z(N(u)),1,0)\bigr),\qquad
 M(n):=j\bigl(j(J(n,0;u.j(u)))\bigr).
\]
Finite computations give
\[
 j(0)=1,\qquad j(1)=2,\qquad j(2)=0.
\]
The range of $j$ is contained in $\{0,1,2\}$: its zero branch is
one, and its successor branch is zero or two according to $N(u)$.
Consequently the definitions and these three computations give
\[
 M(0)=0,\quad M(Sn)=j(M(n)),\quad M(S^3n)=M(n),
 \qquad M(n)\in\{0,1,2\}.
 \tag{Three-cycle}\label{js:three-cycle}
\]
Here $M(0)=j(j(1))=0$, and the successor equation follows by
associativity of composing the fixed term $j$ with itself.  The period
three follows from its three possible values.  No numerical division
algorithm has been assumed.

\subsection{The shifted ramp and factor map}

Set
\[
 G(n):=J(n,0;u.u\oplus M(S^2u)).
\]
Its first values are
\[
 G(0)=1,\ G(1)=2,\ G(2)=4,\ G(3)=5,\
 G(4)=7,\ G(5)=8,\ G(6)=10.
\]

\begin{lemma}[Ramp identities]\label{js:ramp}
The preceding theory proves
\[
 M(G(n))=S(N(n)),\qquad G(S^2n)=S^3G(n),
 \tag{Ramp}\label{js:ramp-identities}
\]
and admits the following two finite-case calculations:
\[
 \begin{array}{c|cc}
 &N(n)=0&N(n)=1\\ \hline
 G(Sn)&S G(n)&S^2G(n)\\
 G(S^3n)&S^4G(n)&S^5G(n).
 \end{array}
 \tag{Ramp cases}\label{js:ramp-cases}
\]
\end{lemma}
\begin{proof}
Induct on the single equation $M(G(n))=S(N(n))$.
Its base is $M(1)=1$.  If $N(n)=0$, the induction equation gives
$M(G(n))=1$, hence $M(S^2G(n))=0$ by the three-cycle.  The next
state is $G(n)\oplus0=S G(n)$, whose $M$-value is two, as
required.  If $N(n)=1$, the current $M$-value is two and
$M(S^2G(n))=1$.  The next state is
$G(n)\oplus1=S^2G(n)$, whose $M$-value is one.  This proves
the invariant and the first line of the case table.  Applying that line
twice, with the parity cycle, gives the two-step equation; applying it
three times gives the last line.
\end{proof}

\begin{lemma}[Affine identities]\label{js:affine-lemmas}
Throughout the entire numerical domain,
\[
 N(D(a))=1,\qquad N(T(a))=N(a),\qquad
 S G(D(a))=T(a),
 \tag{Affine parity and ramp}\label{js:affine-parity}
\]
\[
 G(T(a))=T(G(a))\oplus N(a).
 \tag{Ramp scaling}\label{js:ramp-scaling}
\]
\end{lemma}
\begin{proof}
For $N(D(a))=1$, use EI, $D(0)=3$, and the two-step parity
period.  For $N(T(a))=N(a)$ use EI, $T(0)=6$, and
$T(Sa)=S^3T(a)$; three parity steps toggle both sides once.

For $S G(D(a))=T(a)$ use EI on that literal equation.  The base
is $S G(3)=6$.  The successor step uses
$D(Sa)=S^2D(a)$ and $G(S^2v)=S^3G(v)$, so it applies three
successors to the previous equation, just as $T$ does.

For \eqref{js:ramp-scaling}, the base is
$G(6)=10=T(1)\oplus0$.  Retain its equation as the induction
hypothesis, and split $N(a)$.  In the zero case, $N(T(a))=0$,
so the new left side is $S^4G(T(a))$.  The old right side is
$T(G(a))\oplus0=S T(G(a))$, while the new right side is
\[
 T(SG(a))\oplus1=S^5T(G(a)).
\]
Thus both sides have changed by four successors.  In the one case,
the new left side is $S^5G(T(a))$.  The old right side is
$T(G(a))\oplus1=S^2T(G(a))$, while the new right side is
\[
 T(S^2G(a))\oplus0=S^7T(G(a)).
\]
Thus both sides have changed by five successors.  These calculations
use only \eqref{js:oplus}, \eqref{js:affine-shifts}, and the ramp
case table; no successor has been cancelled.
\end{proof}

Define the factor map
\[
 H(u):=B(Z(N(u)),G(u),u).
\]
Its two cases are
\[
 N(u)=0\ \Longrightarrow\ H(u)=u,\qquad
 N(u)=1\ \Longrightarrow\ H(u)=S G(u).
\]

\paragraph{Why these shifts are useful.}
For intuition on standard natural numbers only, $G(u)=\lfloor3u/2\rfloor+1$.
Thus $H$ fixes even inputs and sends an odd $u$ to $3(u+1)/2$.
After shifting the state by three, an active step multiplies it by $3/2$.
The affine maps $D(a)=2a+3$ and $T(a)=3a+6$ are precisely doubling
and tripling in this shifted coordinate: $D(a)+3=2(a+3)$ and
$T(a)+3=3(a+3)$. Starting at two makes the unshifted base even and
fixed by $H$, while its shifted value is five. The later seed sequence
therefore has standard values $5\cdot2^x-3$. This motivates a comparator
that tests when its powers of two have been exhausted. The proof below
uses the displayed $J$ terms and literal equations; division, exponentiation,
and this standard-value description are not prerequisites.

\begin{lemma}[The factor-map equations]\label{js:factor}
The following global equations hold:
\[
 H(2)=2,\qquad H(D(a))=T(a),\qquad H(T(a))=T(H(a)).
 \tag{Factor}\label{js:factor-equations}
\]
\end{lemma}
\begin{proof}
The first equation follows from $N(2)=0$.  Since $N(D(a))=1$,
the second is $S G(D(a))=T(a)$.
For the last equation split $N(a)$, noting that $N(T(a))=N(a)$.
In the zero case both sides equal $T(a)$.  In the one case,
\eqref{js:ramp-scaling} gives
\[
 H(T(a))=S G(T(a))
 =S\bigl(T(G(a))\oplus1\bigr)
 =S^3T(G(a))=T(SG(a))=T(H(a)).
\]
\end{proof}

\subsection{A comparator without stripping an initializer}

Put
\[
 A(x):=J(x,0;u.R(u)),\qquad E(x):=S A(x),
\]
\[
 F(x,y):=J(y,A(x);u.H(u)),\qquad
 L(x,y):=Z(N(F(x,y))).
\]
By direct computation and shifted algebra,
\[
 E(0)=2,\qquad E(Sx)=D(E(x)),\qquad F(x,0)=E(x).
\]
The definition of $F$ is the decisive initializer alignment:
its successor-seeded initial value is exactly $E(x)=S A(x)$.
There is no hidden use of predecessor to initialize it at $E(x)$.

\begin{lemma}[Comparator equations]\label{js:comparator}
The term $L$ is zero or one and satisfies
\[
 L(0,y)=1,\qquad L(Sx,0)=0,\qquad
 L(Sx,Sy)=L(x,y).
 \tag{Comparator}\label{js:comparator-equations}
\]
\end{lemma}
\begin{proof}
Its range follows from the range of $Z$.  EI on $y$ proves
$F(0,y)=2$, since $E(0)=2$ and $H(2)=2$; hence $L(0,y)=1$.
Also $F(Sx,0)=E(Sx)=D(E(x))$ has parity one, giving
$L(Sx,0)=0$.

For the remaining equation first prove, by EI on $y$ with $x$ fixed,
\[
 F(Sx,Sy)=T(F(x,y)).
\]
At zero, the left side is
$H(E(Sx))=H(D(E(x)))=T(E(x))$.  In the successor step both
iterates apply $H$, and $H\circ T=T\circ H$ rewrites the result
into $T(F(x,Sy))$.  Applying $Z\circ N$ and using
$N(T(v))=N(v)$ gives the desired comparator equation.

This proof avoids cancellation in an equation between two successor
initializers.  Such cancellation would be unavailable at this stage.
\end{proof}

\subsection{Saturation and a finite correction give true predecessor}

Define conditional increment by
\[
 K(c,u):=B(c,u,u),\qquad
 B_m(u):=K(Z(L(Su,m)),u),\qquad
 V(t,m):=J(t,0;u.B_m(u)).
\]
Thus $K(0,u)=Su$ and $K(Sv,u)=u$.  The counter $m$ is a passive
parameter of the iteration defining $V$.

\begin{lemma}[Shifted saturation]\label{js:saturation}
The theory proves
\[
 V(St,S^2k)=S V(t,Sk),\qquad
 V(n,Sn)=Sn,\qquad
 V(Sn,Sn)=V(S^2n,Sn)=Sn.
 \tag{Saturation}\label{js:saturation-equations}
\]
\end{lemma}
\begin{proof}
Numerical cases on $c$ give $K(c,Su)=S K(c,u)$.  Consequently
the comparator equations give
\[
 B_{Sm}(Su)=S B_m(u),\qquad B_{S^2k}(1)=2.
\]
EI on $t$, with $k$ fixed, now proves the first displayed equation.
At zero its two sides equal two.  At a successor, apply the defining
update to the induction equation and then use the conjugacy
$B_{S^2k}(Su)=S B_{Sk}(u)$.

EI gives $L(Sn,n)=0$ directly from the comparator equations.
A separate EI proves $V(n,Sn)=Sn$: the base is $V(0,1)=1$,
and the successor step uses the previously proved uniform shift at
$t=k=n$.  This is substitution into a completed lemma, not a change
of a passive parameter inside an active induction hypothesis.
Finally
\[
 B_{Sn}(Sn)=K(Z(L(S^2n,Sn)),Sn)
           =K(Z(L(Sn,n)),Sn)=K(1,Sn)=Sn.
\]
One and then two more iterator steps give the last two equations.
\end{proof}

Put
\[
 Q(m):=J\bigl(m,0;u.K(Z(L(S^2u,m)),u)\bigr).
\]
This term is a capped predecessor rather than the desired predecessor:
\[
 Q(0)=Q(1)=1,\qquad Q(S^2k)=Sk.
 \tag{Capped predecessor}\label{js:capped}
\]
The first two equations are direct computations.  For the last,
\[
 L(S^2u,S^2k)=L(u,k)=L(Su,Sk)
\]
identifies its update with $B_{Sk}(u)$.  The derived iterator uniqueness
lemma, proved at an independent counter $t$ with $k$ fixed, gives
\[
 J(t,0;u.K(Z(L(S^2u,S^2k)),u))=V(t,Sk).
\]
Substitute $t=S^2k$ and apply \eqref{js:saturation-equations}.

Now define one fixed correction term
\[
 \boxed{\quad
 P(m):=B\bigl(Z(L(3,m)),\ Q(Q(m)),\ L(2,m)\bigr).
 \quad}\tag{Predecessor}\label{js:predecessor}
\]
There is no arbitrary selector hidden in this expression.  Its outer
operation is the original partial selector $B$, whose zero branch adds
exactly one successor.

\begin{theorem}[Predecessor and initializer recovery]\label{js:recovery}
The two-schema successor-seeded scalar theory proves
\[
 P(0)=0,\qquad P(Sn)=n.
\]
For every fixed update $h$, the finite term
\[
 I(n,a;v.h):=
 P\bigl(J(n,a;u.S(h[Pu/v]))\bigr)
 \tag{Initializer compiler}\label{js:initializer}
\]
satisfies
\[
 I(0,a;v.h)=a,\qquad I(Sn,a;v.h)=h[I(n,a;v.h)/v].
\]
\end{theorem}
\begin{proof}
The comparator equations give the three direct values
\[
 P(0)=0,\qquad P(1)=0,\qquad P(2)=1.
\]
Indeed, at these inputs $L(3,m)=0$, so the outer $B$ selects
$L(2,m)$, whose three values are zero, zero, one.
At $m=S^3k$, one has $L(3,m)=1$.  Thus the zero-control branch
of $B$, together with \eqref{js:capped}, gives
\[
 P(S^3k)=S Q(Q(S^3k))=S Q(S^2k)=S^2k.
\]
Numerical cases twice on $n$ now give $P(Sn)=n$ from these displayed
equations.  Alternatively, two literal-equation inductions, first on
$P(S^2n)=Sn$ and then on $P(Sn)=n$, give the same conclusion.
No earlier cancellation or injectivity was used.  Successor injectivity
can now be derived by applying $P$ to an equality of successors.

For the initializer compiler, J0 and the proved predecessor equation
give the base $P(Sa)=a$.  At a successor, JS exposes the term
$S(h[P(J(n,a;u.S(h[Pu/v])))/v])$.  The outer $P$ removes its
displayed successor and leaves exactly $h[I(n,a;v.h)/v]$.
There is no invariant about the internal state to prove: the displayed
successor in the update is sufficient.
\end{proof}

\subsection{Consequences and the bi-interpretation boundary}

The construction supplies the full initialized scalar basis $0,S,I$.
In particular the unrestricted selector is now the derived term
$C(n,a,b):=I(n,a;u.b)$, with its ordinary endpoints $a$ and $b$.
The supplied scalar arithmetic development can therefore be used from
this point onward to obtain actual-equality reflection, induction on
arbitrary quantifier-free formulas, pairing and projections, and the
compiler for every fixed general primitive-recursion description.
Those later arguments are not premises of this bootstrap.

In ordinary iPRA, each fixed $J$ description is interpreted by primitive
recursion with seed $Sa$ and its displayed scalar update.  Both
computation schemas, coverage, separation, and literal-equation induction
are then provable.  Conversely \eqref{js:initializer}, followed by the
existing scalar compiler, interprets every fixed iPRA operation.
For each fixed description, literal-equation induction proves uniqueness
of its defining recurrence after earlier subdescriptions have been
identified.  Applied to both composite translations, this identifies
each numerical operation with its double translation throughout the
entire original domain.  Both domain formulas are $x=x$ and both
numerical isomorphism graphs are $x=y$.

For the ordinary presentations, full intuitionistic logic is preserved.
For the native paired-binder presentation, its ordinary logical macros
and the formula return-translation proof must still be included; they
apply once the derived selector $C$ is available.  This establishes
bi-interpretability only when that logical bridge and the scalar
arithmetic completion accompany the present argument.  It does not
replace them by agreement merely on the standard natural numbers.

The arithmetic advance is completion of the predecessor obligation
left open for this exact candidate in the shared-continuation
report~\cite{shared}: the affine factor-map alignment, capped
predecessor, and finite correction \eqref{js:predecessor} supply its
missing proof.  The candidate grammar and initializer compiler are not
new.  Finite cycles, a $2/3$ factor map, and saturation are adapted from
the supplied ordinary scalar predecessor construction.  No global
literature priority is asserted.

The now-completed basis has two primitive numerical constructors and
two computation schemas, with a single update variable.  A smaller
complete proof checker remains a separate implementation claim, which
must charge the actual decoder, scope checks, logical rules, and driver.


\section{Recovering full first-order intuitionistic logic}
\label{sec:choice-introduction}
After Theorem~\ref{js:recovery}, the derived initialized iterator $I$ and
the genuine selector $C(n,a,b):=I(n,a;u.b)$ are available. Consequently
$C(0,a,b)=a$ and $C(Sn,a,b)=b$ on the whole domain. The following
logical construction is inherited from the supplied native-binder
developments\cite{generated,elimination}; its prerequisites have now
been obtained without circularity.

\subsection{Gates, padding, and formula selection}
First eliminate each $\Q_yB$ by its interderivable representation
$\D_{u,y}B$ with unused $u$. From $\Q$ to $\D$ use DI-case and QE;
from $\D$ to $\Q$ use DE and the supplied Hyp. Hence it suffices to
define the following operations on Eq/$\D$ formulas.

For a fresh control $v$, define a gate by external recursion on formulas:
\[
 G_v(s=t):=(C(v,0,s)=C(v,0,t)),\qquad
 G_v(\D_{x,y}A):=\D_{x,y}G_vA.
\]
It satisfies $\vdash G_0A$ and $G_{Sv}A\dashv\vdash A$.
Uniform interderivability lifts through $\D$: open the original binder,
convert the new base to the old base, use its closure at a successor,
convert the result, and introduce the new binder by DI-case.
This proves the gate identities by external structural induction.
It also proves the padding equivalence
\[
 A\dashv\vdash\D_{u,y}G_yA\qquad(u,y\notin\FV(A)).
\]
The forward direction is DI-case. For the reverse, DE, the zero-gate
theorem, and QE at $1$ recover $A$.

Define the formula selector by
\[
 \mathcal M_v(s=t,r=q):=(C(v,s,r)=C(v,t,q)),
\]
\[
\begin{aligned}
 \mathcal M_v(\D_{x,y}A,\D_{x,y}B)&:=\D_{x,y}\mathcal M_v(A,B),\\
 \mathcal M_v(E,\D_{x,y}B)&:=\D_{x,y}\mathcal M_v(G_yE,B),\\
 \mathcal M_v(\D_{x,y}A,E)&:=\D_{x,y}\mathcal M_v(A,G_yE).
\end{aligned}
\]
Here $E$ is an equation and bound coordinates are freshly aligned.
Every recursive call decreases maximum binder depth. Padding and binder
congruence prove $\mathcal M_0(A,B)\dashv\vdash A$ and
$\mathcal M_{Sv}(A,B)\dashv\vdash B$, including different input depths.

\subsection{Explicit logical macros and their rules}
Put $\delta(v)=C(v,1,0)$ and $p(z)=I(z,0;u.\delta(u))$.
Its fixed computations give $p(1)=1,p(2)=0$. The entire macro list is
\[
\begin{aligned}
 A\Rightarrow_*B&:=\D_{u,z}\mathcal M_z(A,B),\\
 A\wedge_*B&:=\D_{u,z}G_z\mathcal M_{p(z)}(A,B),\\
 \exists_*xA(x)&:=\D_{x,y}G_yA(x),\\
 A\vee_*B&:=\exists_*x\mathcal M_x(A,B),\\
 \forall_*xA(x)&:=A(0)\wedge_*\D_{u,z}G_zA(z),\\
 \bot_*&:=(0=1).
\end{aligned}
\]
All auxiliary variables are fresh and $u$ is unused. Implication
introduction is DI-case, using the selector's zero endpoint as its local
assumption; implication elimination opens the binder and specializes at
$1$. Existential introduction uses the supplied witness and the gate;
existential elimination opens its witness and specializes the closure
using the zero-gate theorem.

Conjunction introduction first uses NC to obtain
$\mathcal M_v(A,B)$ uniformly from the two supplied proofs, then
substitutes $v=p(Sz)$ for the successor matrix. Its projections specialize
at $1$ and $2$. Disjunction introduction uses selector witness $0$ or
$1$. For elimination, form the two branch implications into the common
conclusion; numerical cases give the implication from an arbitrary
selector, and existential elimination supplies that selector.

These derived rules and numerical cases prove
$x=0\vee_*\exists_*y(x=Sy)$. Universal introduction supplies its zero
and successor parts from a proof at a fresh parameter. Universal
elimination uses that coverage: its zero branch is the first conjunct,
and its successor branch is obtained from the gated closure. Under
$0=1$, the selector equations give every numerical equation; external structural
induction on formulas using DI-case gives every native formula; a
$\Q$ formula uses its unused-$\D$ representation. Thus
falsity elimination is available.

For explicitness, the recovered ordinary inference inventory is Hyp,
equality reflexivity and replacement, falsity elimination,
$\to$I/$\to$E, $\wedge$I and both projections, both $\vee$I rules
and $\vee$E, $\forall$I/$\forall$E, and $\exists$I/$\exists$E.
The last four are
\[
 \frac{\Gamma\vdash A(z)}{\Gamma\vdash\forall xA(x)},\quad
 \frac{\Gamma\vdash\forall xA(x)}{\Gamma\vdash A(t)},\quad
 \frac{\Gamma\vdash A(t)}{\Gamma\vdash\exists xA(x)},\quad
 \frac{\Gamma\vdash\exists xA(x)\quad\Gamma,A(z)\vdash F}
      {\Gamma\vdash F}.
\]
Universal introduction excludes $z$ from $\Gamma$ and passive data;
existential elimination also excludes it from $F$. The propositional
rules have their usual reusable-context form: introduce an implication
by discharging its antecedent, eliminate it with a proof of that
antecedent; form a conjunction from its two proofs and project either;
introduce either disjunct and eliminate a disjunction using two local
proofs of the same conclusion. No additional classical principle is
silently included.

\section{The two interpretations and their composites}
\label{sec:interpretations}
Let $T_J$ be the ordinary first-order realization: intuitionistic equality
logic, J0/JS, literal-equation induction, numerical coverage, and
$(0=1)\to\bot$. The explicit falsity bridge is essential; the positive
equations alone admit a singleton structure. The native theory interprets
$T_J$ by the preceding macros, equation induction, coverage, and explosion.
Conversely the ordinary meanings of $\Q,\D$ validate every native rule:
DI-case uses coverage, and DI-progress uses induction on its literal
equation. The supply presentation is equivalent by substitution and cut.

Appendix~\ref{sec:eic-completion}, instantiated with the derived $P,C,I$,
constructs a characteristic of actual equality and translates every
quantifier-free formula into a characteristic equation. It therefore
recovers all quantifier-free induction. Appendix~\ref{sec:iterator-coding}
constructs pairing and its projections by bounded scalar iteration and
compiles each fixed general PR description. Pairing is only a state-coding
tool; the interpretation domain remains the whole original domain.

\begin{theorem}[Full-domain bi-interpretability]\label{thm:main}
The stated native theory, its computation-supply presentation, and $T_J$
are faithful presentations of a theory bi-interpretable and equiconsistent
with full first-order $\ipra$. Both numerical domain formulas are $x=x$;
both numerical composite isomorphisms have graph $x=y$.
\end{theorem}
\begin{proof}
In $\ipra$, interpret each fixed $J$ body by primitive recursion with
base $Sa$ and step $h$, ignoring the recursion counter in the update.
Its two equations and ordinary boundary and induction rules are provable.
In the other direction, the initializer compiler of
Theorem~\ref{js:recovery}, followed by the scalar PR compiler, gives a
finite term for each fixed target description. Equality reflection
validates the constructor laws and every target induction instance.

For the numerical composites, induct externally on each finite
description. Once earlier subdescriptions are identified, the original
operation and its double translation have the same base and update.
Internal induction on the single equation between their values proves
equality at every counter. Parameters remain fixed during this induction.
For the original PR descriptions use the corresponding uniqueness proof
with their counter argument. This proves internal operation preservation
by the identity relation, not merely equality of standard numerical
outputs. The same argument gives coherence when substituting parameters
changes a staged update description.

For the formula composites, the six displayed ordinary logical macros
have their intended meanings by their endpoint and padding proofs.
The nontrivial native return translations are
\[
 \Q_yB\dashv\vdash B(0)\Rightarrow_*\forall_*yB(y),
\]
\[
 \D_{x,y}B\dashv\vdash
 \exists_*x\bigl(B(x,0)\Rightarrow_*\forall_*yB(x,y)\bigr).
\]
Forward, DE opens the witness and supplies its closure. Under the local
base, QE supplies the body at a fresh parameter; introduce the universal,
implication, and existential. Backward, eliminate the existential and use
its implication and universal to supply DI-case's successor premise under
the retained base. The $\Q$ case follows directly or from its unused
witness representation. External induction on every finite native and
ordinary formula gives both formula composites, in arbitrary contexts.

Finally, ordinary falsity and native $0=1$ are contradictions in both
directions. The effective finite proof translations therefore give
equiconsistency. No theory is asserted to prove its own consistency.
\end{proof}

\clearpage
\section{Complete executable checkers and their correspondence}
\label{sec:implementation}

The numerical completion does not by itself shrink a validator. The
programs below retain the generated-assumption architecture of the
supplied draft and change its arithmetic base equation from $I(0,a)=a$
to $J(0,a)=Sa$. The byte profile also replaces one conditional statement
by an equivalent conditional expression. The AST profile uses the
exception-based opener and the explicitly decoded iterable aliases
described below. Every entire executable module is charged in the
comparison, including the JSON input and verdict driver.

\Needspace{28\baselineskip}
\subsection{The complete byte profile}
\begin{lstlisting}[language=Python]
from json import load
k = 0
i = lambda a, t=None, d=0: [i(b, t, d + b' \x02\x01\x04'[len(a)]) for b in a] if isinstance(a, list) else a if k <= a < d else d % 2 and i(t, d=1)

def c(a, g):
    global k
    if i(a) in g:
        return
    m, s, u, t = q()
    if m:
        c(m, g)
        a = i(t, s) if m in ([s, u], not s and [t]) and i(t, u) == a else a
        h = m[2:]
    elif t:
        h = a = i(a[2], t, 2)
    else:
        n, t, u = s
        h = [s, i(u, i([*n, t, u], d=1), 1) if n else [t]]
    k -= 1
    c(i(a, [k]), g + [i(h, k if m or len(a) % 2 == m else 0)])
try:
    q = load(open(0)).pop
    c(q(), [])
except:
    print('REJECT')
    raise
print('ACCEPT')
\end{lstlisting}
The imported JSON parser supplies data only. The submitted data is never
evaluated as Python source. No arithmetic oracle, external theorem table,
proof producer, helper validator, or unmeasured module is imported.
The arithmetic change is the final \texttt{[t]} in the assignment of
\texttt{h}; that is the raw successor of the initializer.

\Needspace{30\baselineskip}
\subsection{The complete AST profile}
\begin{lstlisting}[language=Python]
from json import load
k = 0

def i(a, t=None, d=0):
    try:
        return [i(b, t, d + b' \x02\x01\x04'[len(a)]) for b in a]
    except:
        return a if k <= a < d else d % 2 and i(t, d=1)

def c(a, g):
    global k
    if i(a) in g:
        return
    m, s, u, t = q()
    if m:
        c(m, g)
        if m in ([s, u], not s and [t]) and i(t, u) == a:
            a = i(t, s)
        h = m[2:]
    elif t:
        h = a = i(a[2], t, 2)
    else:
        n, t, u = s
        h = [s, i(u, i([*n, t, u], d=1), 1) if n else [t]]
    k -= 1
    c(i(a, [k]), g + [i(h, k if m or len(a) % 2 == m else 0)])
try:
    q = load(open(0)).pop
    c(q(), [])
except:
    print('REJECT')
    raise
print('ACCEPT')
\end{lstlisting}
The two checkers have the same canonical mathematical certificates.
The latter additionally accepts the safe non-array aliases specified
below. The different implementations improve different coordinates;
their costs must not be combined into an imaginary single program.

The additional complete file \file{code/checker_ast_total.py} replaces
\texttt{for b in a} by \texttt{for b in a.copy()} and binds the input
pop method with the conventional assignment expression
\texttt{c((q := load(open(0)).pop)(), [])}. It measures $287/314/711$.
The input is still loaded once and its final goal is popped before the
checker runs. Strings lack \texttt{copy} and immediately reject; a
nonempty object's string keys likewise reject. Only the empty object
adds an alias of $[]$. Opening therefore terminates structurally on
finite input trees, including malformed ones. Each non-hypothesis proof
invocation consumes a packet. This alternative improves the old AST
frontier by one node without the malformed-string tradeoff of the
313-node version. Its source is included in full in the archive.

\subsection{Public protocol and exact sorted decoding}

The input is one JSON array. Its last item is the closed requested
formula. Earlier items are four-field instructions $[m,s,u,t]$, consumed
from right to left by \texttt{pop}. A major subproof is consumed before
its continuation. A hypothesis closes without consuming an instruction.
An unread prefix is inert. Successful execution prints
\texttt{ACCEPT} and exits with status zero; rejection prints
\texttt{REJECT} and exits nonzero, possibly with a traceback on standard
error. Resource limits can reject sufficiently large valid certificates;
mathematical completeness is for unbounded resources.

\begin{center}
\begin{tabular}{cllc}
\toprule
Arity & Formula & Numerical term & Depth increase\\
\midrule
0 & truth $0=0$ & zero & no children\\
1 & conditional closure $\Q$ & derived successor & 2\\
2 & equality & annotated zero & 1\\
3 & paired closure $\D$ & successor-seeded $J$ & 4\\
\bottomrule
\end{tabular}
\end{center}
The unary numerical node is an abbreviation for $J(0,a;u.0)$.
It is a charged serialization convenience, not an additional primitive
computation schema. Numerical binary nodes denote zero but their
annotations are still traversed. The first two fields of a ternary
formula are likewise checked annotations.

For a protected numerical environment $\rho$, nearest slot first, the
precise nontrivial decoding clauses are
\[
\begin{aligned}
 T_\rho([p])&=S T_{0,\rho}(p),&
 F_\rho([B])&=\Q_y F_{y,\rho}(B),\\
 T_\rho([p,q])&=0,&
 F_\rho([s,t])&=(T_\rho(s)=T_\rho(t)),\\
 T_\rho([n,a,h])&=
 J(T_{0,0,\rho}(n),T_{0,0,\rho}(a);v.T_{v,v,\rho}(h)),\\
 F_\rho([P,U,B])&=\D_{x,y}F_{y,x,\rho}(B).
\end{aligned}
\]
Empty arrays decode as zero at numerical sort and truth at formula sort.
The two update slots alias the \emph{same} scalar state; they are not
separate changing inputs. The two slots of $\D$ are distinct numerical
binders. The extra slots under a successor and under an iterator's
counter and initializer are dummy zeros.

At formula depth $2q$ and numerical depth $2q+1$, there are $q$ protected
slots. Positive coordinates $2j-1$ and $2j$ name the same $j$th slot;
zero denotes numerical zero. Negative coordinates name allocated rigid
parameters. The global lower bound $k$ begins at zero and decreases at
every non-hypothesis continuation, including computation and neutral
endpoint continuations. A major subproof may allocate names before the
outer continuation allocates its own name.

The opener $i(a,t,d)$ retains a scalar in $[k,d)$. An exposed numerical
coordinate is replaced by the supplied image, itself checked at numerical
root, depth one. Exposed formula coordinates denote truth. A missing or
invalid \emph{used} numerical image rejects. This numerical-root check is
what prevents capture by a surrounding template. The requested goal is
checked with no supplied image before any continuation, so meaningful
future numerical parameters cannot occur in it.

Canonical data uses arrays and integers. Retained finite fractional
coordinates alias their integer floors, and Boolean values have their
zero/one meanings. An exceptional scalar can survive only by being
replaced with a successfully checked image, or as an exposed formula
truth. These conventions do not add numerical operations or axioms.

\subsection{The AST profile's additional local aliases}

The AST opener traverses its argument directly inside \texttt{try}.
On a successful used opening, the only additional iterable possibilities
are:
\[
 \texttt{""}\rightsquigarrow[],\qquad
 \texttt{\{\}}\rightsquigarrow[],\qquad
 \texttt{\{"":value\}}\rightsquigarrow[[]].
\]
The value of the sole dictionary entry is ignored. A nonempty string
recurses on a one-character string without reaching a leaf, so cannot
finish successfully. A dictionary with a nonempty key therefore cannot
complete a used opening either. The empty-key singleton has exactly the
length, truth value, and traversed child of a unary node. All three
normalizations commute with scope opening. For the first two aliases the
depth increment is immaterial because there are no children.

Actual CPython rejects the nonempty-string case through its finite
recursion limit. An ideal interpreter with no recursion bound would
diverge on that malformed input. Thus the AST reduction has this explicit
raw-input tradeoff; the byte profile retains list dispatch and avoids it.
The soundness statement concerns every finite successful run, and the
completeness statement promises canonical array/integer certificates.

An empty-key dictionary cannot successfully be a non-hypothesis major:
the subsequent slice \texttt{m[2:]} fails. As a unary successor counter,
its expansion supplies the empty-string zero predecessor, whose numerical
opening normalizes to zero. Thus it gives the same computation as the
corresponding array. These are local statements about successfully opened
and inspected data, not a claim that every ignored field in a raw
certificate has a global canonical decoding. Arbitrary unused fields
and unread prefixes remain semantically inert.

\subsection{Every instruction branch}

\paragraph{A truthy major.}
First recursively prove $m$; a failed major is never ignored. If
$m=[s,u]$ and the target opening of the formula motive $t$ at $u$
matches the requested goal, continue at its opening at $s$. This is
equality transport. If $s$ is falsy and $m=[t]$, the analogous match is
QE, with its zero-instance minor. In either case the extra context entry
is truth, the empty slice $m[2:]$. An unmatched binary or unary major
retains the requested goal and adds only that truth. It does not supply
the alleged endpoint equality.

A ternary major cannot match either endpoint candidate. Its slice is
$[B]$; opening at the newly allocated name supplies the $\Q$ assumption
of DE, while the requested conclusion is retained. The fresh name is
absent from the already checked conclusion and context.

\paragraph{A falsy major and a truthy original fourth field.}
This is paired introduction with witness $t$. The opener takes the goal's
third field and leaves one induction coordinate protected. Marker $m=[]$
chooses its zero-instance assumption. Marker $m=0$ chooses the current
instance only when the opened matrix is a literal equation; otherwise
it also chooses the zero instance. The new goal is the successor
instance. The test \texttt{len(a) \% 2 == m} is precisely this distinction
on successfully decoded matrices. It never permits progress introduction
on a $\Q$ or $\D$ matrix.

A zero witness must have a truthy encoding in this branch. In the new
basis one can use the term $J(1,0;u.0)$, which is provably zero. The old
initialized-iterator witness $[0,0,0]$ now denotes one and cannot be reused
as zero. The supplied producer makes the required change. Routine
transport along a proved equality handles the selected zero encoding.

\paragraph{A falsy major and falsy original fourth field.}
The source $s$ must unpack as $[n,a,h]$. The packet's original third
field is unused. With a false counter the added raw equation is
$[s,[a]]$, expressing the successor-seeded base. With a true counter
the inner opening of $[*n,a,h]$ both forces the counter to be unary and
validates the predecessor iterator before substitution into the update.
This eager check is essential even when that update ignores its state.

Every branch then decrements $k$, opens the goal with the fresh successor
$[k]$, adds the appropriately opened assumption, and uses the same
recursive continuation. The major and parity tests select a fresh-name
assumption for DE/progress and a zero assumption for cases; computation
can generate a local equation containing the fresh name. The admissible
known-instance supply rule accounts for precisely that last possibility.

\subsection{Why generated J equations are sound}

\begin{lemma}[Structural zero extension]\label{impl:zero}
Read a fixed numerical tree under additional outer slots assigned zero.
Alternatively expose those slots to a fresh rigid parameter $w$ and then
substitute $w=0$. Whenever the indicated openings succeed, the decoded
numerical terms agree in the theory.
\end{lemma}
\begin{proof}
Induct externally on the finite tree. Internal binder slots and older
rigid parameters are unchanged. An added zero slot, an exposed coordinate,
or an alias of $w$ becomes zero. Successor adds a dummy zero on both
sides. A checked numerical binary node denotes zero on both sides.
At a ternary node its counter and initializer add two zero slots and its
update adds two aliases of one symbolic state. Apply the induction
hypothesis to all three children. Their fixed $J$ descriptions agree by
substitution and congruence. The local iterable aliases in the AST profile
normalize before the same argument. This proof does not enumerate the
standard natural numbers.
\end{proof}

\begin{lemma}[The generated equation has a provable zero instance]
\label{impl:supply}
If a computation instruction successfully adds $H(w)$, where $w$ is its
fresh parameter, then the native theory proves $H(0)$.
\end{lemma}
\begin{proof}
At a false counter the left side is $J(0,a;v.h)$. Its initializer is
read under two dummy zero slots. The initializer on the right side is
under one dummy zero, inside its successor, and may expose additional
coordinates to $w$. After $w=0$, Lemma~\ref{impl:zero} identifies the
two initializers. The equation is consequently J0. The particular update
does not affect this base value.

At a true counter the eager predecessor check forces a unary $n=[p]$,
or its documented empty-key alias. An array counter with two or more
fields makes a predecessor tree of arity at least four, which the depth
table rejects. A scalar cannot be expanded. The predecessor tree, its
initializer, and its update are numerically checked before fresh-name
allocation. Thus an out-of-scope coordinate cannot masquerade as its
recursive state. The right side substitutes exactly that checked
predecessor iterator into the update. The extra dummy zero under the
source successor changes no surviving meaningful coordinate, by the
same structural argument. Its zero instance is JS.
\end{proof}

The original conclusion and retained context do not contain $w$.
Substitute zero for $w$ in the continuation derivation and cut the
proved $H(0)$. The local equation need not be valid for every value of
$w$; its known zero instance is sufficient. This proof is the inherited
generated-assumption argument with its base adapted to the new seed.

\begin{theorem}[Checker correspondence]\label{impl:correspondence}
Every successful finite execution of each delivered complete checker has a
decoded closed conclusion provable in the stated native theory.
Conversely, every closed native theorem has a finite canonical certificate
accepted by each checker, with unbounded computational resources.
\end{theorem}
\begin{proof}
For soundness, induct on the finite successful checking trace. Hyp closes
a retained formula. A matched transport or QE command is sound by
capture-free opening. An unmatched unary or binary major is neutral,
and a ternary major performs DE. The DI command supplies exactly one of
the two permitted assumptions and maintains its eigenparameter condition.
A computation continuation is admissible by Lemma~\ref{impl:supply} and
known-instance supply.

Every added context entry is an array produced by opening. For the
strict list profile a scalar requested formula cannot close by membership
in this list context; its other branches reject or retain it while
consuming another instruction. Hence there is no finite successful scalar
goal trace. For the AST profile, the additional successfully opened
iterables are exactly the local array aliases above. This establishes the
required root sorts without an omitted predicate. The initial scope check,
the eager numerical-image check, and the decreasing floor maintain the
freshness and substitution invariants throughout the trace.
The total-input AST variant has only the empty-object alias, and the
assignment expression preserves the input driver's evaluation order.

For completeness, serialize a finite native derivation in the checker's
depth-first consumption order. Replace computation leaves with supply of
their constant equation followed by Hyp. Use matched endpoints for
transport and QE. Use the designated DI markers and a truthy witness
encoding, with equality expansion when needed. Canonical descriptions put
zero in dummy positions and use allocated rigid names for passive
parameters. Allocate a new label at every common continuation, including
neutral and arithmetic ones, and account for automatic hypothesis closure.
Thus all steps have finite canonical certificates. This is an effective
metatheoretic proof translation, not a claim that the small example
producer implements the entire translation.
\end{proof}

\subsection{Executable evidence and its limits}
The main suite contains 36 distinct fixtures: 33 common certificates
(17 accepted derivations and 16 rejected inputs), plus three intentional
raw-alias boundaries between the profiles. All 216 complete-process
runs passed: each fixture was run on all three complete checkers, both
normally and with Python \texttt{-O}. All 24 targeted mutation controls
passed. They detect removal of major-premise checking, admission of
compound progress, removal of eager recursive-image checking, and
restoration of the incorrect old initialized base equation.

A focused AST suite additionally passed 90 complete-process runs,
covering the 33 common cases and 12 explicit alias controls in both
execution modes. Some cases intentionally overlap the main suite.
The receipts record exact source hashes, expected verdicts, exit codes,
and individual results. Accepted examples include active-state JS,
induction that uses its actual hypothesis, retained assumptions, nested
binders, a nonvacuous paired witness, arbitrary QE specialization,
compound-formula transport, and safe neutral endpoint requests.

The total-input AST alternative passes its own 94 complete-process runs:
the 33 common cases and 14 focused controls, each normally and with
\texttt{-O}. Its receipt verifies the exact frozen source and the
different alias boundary.

The separate numerical diagnostic evaluates the displayed $J$ definitions
and checks 306 selected identities and values. It is a finite diagnostic,
not the whole-domain proof. The latter is the equation-induction argument
in the preceding sections. The named certificate producer demonstrates
nonvacuous logical and computation rules, but does not emit the entire
predecessor derivation as a kernel certificate. Neither finite testing nor
independent assistant review is represented as a machine-checked
metatheorem. The archive records this boundary in \file{CLAIMS.json}.

Run \texttt{python3 reproduce.py} for the supplied measurement and proof
fixtures, and \texttt{python3 build.py} to compile the PDF. The frozen
source archives and their hashes allow the corpus measurements and
textual novelty search to be repeated without relying on a moving branch.


\Needspace{25\baselineskip}
\section{Further research questions}
\begin{enumerate}
\item Can the shifted comparator and capped-predecessor correction be
replaced by a substantially shorter $J$ term? The present proof resolves
existence and whole-domain correctness, not optimal term size.
\item Can the native finite-case lifting argument be formalized first,
followed by the comparator and predecessor, in Lean or Rocq? This would
check the most delicate dependency boundary before formalizing the
inherited logical macros and PR compiler.
\item Can the completed numerical basis remove a mechanism from the
complete checker, producing a single implementation below all three
current size coordinates? The mathematical signature reduction alone
does not achieve that stronger benchmark.
\item What lower bounds can be proved after fixing a concrete checker
language, input grammar, runtime primitives, and required rejection
contract? The observed corpus envelope is not such a theorem.
\item Which other initializer maps $\sigma(a)$, in place of $Sa$, permit
recovery of arbitrary initialization from scalar recursion and equation
induction? A useful criterion must supply an internally proved inverse,
not assume one in the compiler.
\item Can the capped-predecessor correction be isolated as a general
finite-exception repair principle for restricted arithmetic bases? The
present term uses only the original partial selector and the comparator.
\end{enumerate}

\appendix
% Standalone section.  Requires only amsmath and amssymb in the parent document.
% No theorem environments, macros, or bibliography keys are assumed.
% All arithmetic claims in this section are proved here.

\section{From literal-equation induction to full intuitionistic PRA}
\label{sec:eic-completion}\label{app:arithmetic}

This section supplies an unconditional completion theorem. It is adapted
with attribution from the current forward-left draft~\cite{forward}. In the
scalar application used here, every function occurring in the proof is
already a scalar term: predecessor is supplied by
Theorem~\ref{js:recovery}, the selector is
$C(n,a,b)=I(n,a;u.b)$, addition is $x+y=I(y,x;u.Su)$, and
truncated subtraction is $d(x,y)=I(y,x;u.Pu)$. All other functions below
are finite compositions of these. Thus this proof does not presuppose
that arbitrary primitive recursion has already been compiled. The QFI
conclusion first applies to the scalar language; the general PR compiler
comes afterward.  Literal-equation
induction, together with the primitive-recursive defining equations,
predecessor coverage, and separation of $0$ and $1$, proves induction for every
quantifier-free formula.  In particular, the proof does not assume a faithful
numerical equality test, stability or decidability of equality, successor
uniqueness, or any commutation law for an iterator.  We construct the equality
test and prove its reflection property for arbitrary elements of the
numerical domain.

\subsection{The exact hypotheses}

Work in ordinary intuitionistic first-order logic with equality and genuine
falsity $\bot$.  The language has one numerical sort and a symbol for each
fixed, finite primitive-recursive description.  Distinct descriptions are
not identified merely because they compute the same standard function.
Let $\mathsf{PRE}$ contain the defining equations for zero, successor,
projections, composition, and primitive recursion.  In particular, choose
descriptions $P$ and $C$ with equations
\begin{align}
 P(0)&=0, & P(Sx)&=x, \label{eq:eic-predecessor}\\
 C(0,a,b)&=a, & C(Sx,a,b)&=b. \label{eq:eic-selector}
\end{align}
Write $1=S0$.  In addition to $\mathsf{PRE}$, assume
\begin{align}
 \mathsf{Cov}:&\qquad
       \forall x\,\bigl(x=0\lor x=S(Px)\bigr),
       \label{eq:eic-coverage}\\
 \mathsf{Sep}:&\qquad 0=1\longrightarrow\bot,
       \label{eq:eic-separation}
\end{align}
and induction for each \emph{literal equation} $E(n,\vec p)$:
\begin{equation}
 \forall\vec p\,
 \left[
   \left(E(0,\vec p)\land
       \forall n\bigl(E(n,\vec p)\longrightarrow E(Sn,\vec p)\bigr)
   \right)
   \longrightarrow\forall n\,E(n,\vec p)
 \right].
 \label{eq:eic-induction}
\end{equation}
Here the schematic expression $E$ is a single equality of numerical
terms; the displayed universal quantifiers range only over numbers.
Every description and every instance of the schema is finite external
syntax.

The axiom formulation \eqref{eq:eic-induction} and the usual contextual
induction rule are interchangeable.  For example, if
\[
 \Gamma\vdash E(0,\vec p),
 \qquad
 \Gamma,E(z,\vec p)\vdash E(Sz,\vec p),
\]
where $z$ is absent from $\Gamma$ and the passive parameters, discharge the
second premise's equation and universally introduce $z$.  The axiom then
gives $\Gamma\vdash\forall n E(n,\vec p)$, which can be instantiated at any
target term.  Conversely, assume the antecedent of the axiom and apply the
contextual rule with fresh variables.  This uses the ordinary intuitionistic
logical rules and permits arbitrary, including quantified, formulas in
$\Gamma$.

\subsection{Constructor separation and two elementary tests}

Define the sign description and the countdown description by
\begin{equation}
 \operatorname{sg}(u)=C(u,0,1),
 \qquad
 K(u)=C(u,1,Pu).
 \label{eq:eic-sign-countdown}
\end{equation}
Their defining equations imply
\begin{align}
 \operatorname{sg}(0)&=0,
 &\operatorname{sg}(Sx)&=1,\label{eq:eic-sign-computation}\\
 K(0)&=1,
 &K(Sx)&=x.\label{eq:eic-countdown-computation}
\end{align}

\paragraph{Constructor laws.}
Successor is injective: applying $P$ to $Sx=Sy$ gives $x=y$ by
\eqref{eq:eic-predecessor}.  Successor never equals zero.  Indeed,
$Sx=0$ gives $1=0$ by applying $\operatorname{sg}$ and using
\eqref{eq:eic-sign-computation}; symmetry and
\eqref{eq:eic-separation} give $\bot$.  Consequently $1\ne0$, where
$A\ne B$ abbreviates $A=B\to\bot$.  None of these arguments uses
induction.

\paragraph{The sign test.}
For every $u$,
\begin{align}
 &\operatorname{sg}(u)=0\lor\operatorname{sg}(u)=1,
       \label{eq:eic-sign-range}\\
 &\operatorname{sg}(u)=0\quad\longleftrightarrow\quad u=0,
       \label{eq:eic-sign-zero}\\
 &\operatorname{sg}(u)=1\quad\longleftrightarrow\quad u\ne0.
       \label{eq:eic-sign-one}
\end{align}
To prove the range assertion, split \eqref{eq:eic-coverage}.  The branch
$u=0$ gives sign $0$, and the branch $u=S(Pu)$ gives sign $1$.
The implication $u=0\to\operatorname{sg}(u)=0$ is computation.
For its converse, the successor branch of coverage would make
$\operatorname{sg}(u)=0$ imply $1=0$, a contradiction; in that branch
explosion gives $u=0$ as well.  Similarly, if
$\operatorname{sg}(u)=1$ and $u=0$, computation gives $0=1$.
Conversely, $u\ne0$ excludes the zero branch of coverage, and the
successor branch computes the sign as $1$.

Coverage and \eqref{eq:eic-selector} also give, for all $u,v$,
\begin{equation}
 C(u,v,v)=v.
 \label{eq:eic-selector-idempotence}
\end{equation}
Both coverage branches compute to $v$, so this identity again needs no
induction.

\paragraph{The countdown test and the notation $W$.}
For every $u$,
\begin{equation}
 K(u)=0\quad\longleftrightarrow\quad u=1.
 \label{eq:eic-countdown-zero}
\end{equation}
If $u=1=S0$, \eqref{eq:eic-countdown-computation} gives $K(u)=0$.
Conversely, assume $K(u)=0$ and split coverage.  If $u=0$, then
$K(u)=1$, contradicting separation, and explosion gives $u=1$.
If $u=S(Pu)$, computation gives $K(u)=Pu$; hence $Pu=0$ and
$u=S0=1$.

If a presentation uses a description named $W$ with defining equations
\[
 W(0)=1,\qquad W(Sx)=x,
\]
then it proves $K(u)=W(u)$ for every $u$ already from coverage.
In the zero branch both sides are $1$; in the branch $u=S(Pu)$ both
sides are $Pu$.  Thus $K$ and $W$ are provably equal on the whole domain,
even if they are distinct descriptions.  This equality does not appeal to
standard-model agreement or to uniqueness of an arbitrary recurrence.

\subsection{Four elementary uses of equation induction}

Choose descriptions for addition and truncated subtraction, with recursion
on the second argument:
\begin{align}
 x+0&=x, &x+Sy&=S(x+y),\label{eq:eic-addition}\\
 d(x,0)&=x, &d(x,Sy)&=P(d(x,y)).\label{eq:eic-subtraction}
\end{align}
Using these equations and \eqref{eq:eic-induction}, we now derive
\begin{align}
 0+y&=y,\label{eq:eic-left-zero}\\
 (Sx)+y&=S(x+y),\label{eq:eic-left-successor}\\
 d(Sx,Sy)&=d(x,y),\label{eq:eic-cancel-successors}\\
 d(Sa,a)&=1.\label{eq:eic-diagonal}
\end{align}
Each induction below has a single equation as its induction formula.

For \eqref{eq:eic-left-zero}, induct on $y$.  The base is $0+0=0$.
The step is $0+Sy=S(0+y)=Sy$.

For \eqref{eq:eic-left-successor}, hold $x$ fixed and induct on $y$.
At zero both sides are $Sx$.  At a successor the induction hypothesis
gives
\[
 (Sx)+Sy=S\bigl((Sx)+y\bigr)
       =S\bigl(S(x+y)\bigr)=S(x+Sy).
\]
Combining \eqref{eq:eic-left-successor} with the defining equation for
addition gives the useful rearrangement
\begin{equation}
 (Sy)+r=y+Sr.
 \label{eq:eic-shift-successor}
\end{equation}
Indeed, both sides are $S(y+r)$.

For \eqref{eq:eic-cancel-successors}, hold $x$ fixed and induct on $y$.
The base computes as
\[
 d(Sx,S0)=P(d(Sx,0))=P(Sx)=x=d(x,0).
\]
Applying $P$ to the induction equation gives the step:
\[
 d(Sx,S(Sy))=P(d(Sx,Sy))=P(d(x,y))=d(x,Sy).
\]

Finally, induct on $a$ for \eqref{eq:eic-diagonal}.  The base is
$d(S0,0)=S0=1$.  At a successor,
\[
 d(S(Sa),Sa)=d(Sa,a)=1,
\]
using \eqref{eq:eic-cancel-successors} and the induction hypothesis.

\subsection{An equation invariant that reflects actual equality}

The remaining arithmetic ingredient is the following equation, valid for
all $x,y$:
\begin{equation}
 C\bigl(\operatorname{sg}(d(x,y)),\ x,\ y+d(x,y)\bigr)=x.
 \label{eq:eic-residual-invariant}
\end{equation}
Its meaning is that a nonzero subtraction residual reconstructs the
original argument by addition.  The selector expresses this assertion as
one literal equation, so equation induction applies directly.

Hold $x$ fixed and induct on $y$.  At $y=0$, subtraction gives residual
$x$, and \eqref{eq:eic-left-zero} gives $0+x=x$.  The desired equation
is consequently $C(\operatorname{sg}(x),x,x)=x$, which is
\eqref{eq:eic-selector-idempotence}.

For the step, assume \eqref{eq:eic-residual-invariant} at $y$ and
abbreviate
\[
 r=d(x,y),\qquad r'=Pr=d(x,Sy).
\]
These are local term abbreviations, not additional function symbols.
Split coverage for $r'$.

If $r'=0$, then $\operatorname{sg}(r')=0$, and the successor instance
of \eqref{eq:eic-residual-invariant} reduces immediately to $x=x$.

In the other branch, $r'=S(Pr')$, so constructor separation gives
$r'\ne0$.  Split coverage for $r$.  Its zero branch would imply
$r'=P0=0$, a contradiction, from which the desired step equation follows.
In its successor branch,
\[
 r=S(Pr)=Sr',\qquad \operatorname{sg}(r)=1.
\]
The induction equation therefore computes to $y+r=x$.  Meanwhile,
\eqref{eq:eic-shift-successor} gives
\[
 (Sy)+r'=y+Sr'=y+r=x.
\]
Since $r'$ is a successor, $\operatorname{sg}(r')=1$, so this is
exactly the branch selected in the successor instance of
\eqref{eq:eic-residual-invariant}.  This completes the equation induction.
All case distinctions came from coverage; no induction on an implication
or a disjunction was used.

\paragraph{Reflection at residual $1$.}
We have, for arbitrary $a,b$,
\begin{equation}
 d(Sa,b)=1\quad\longleftrightarrow\quad a=b.
 \label{eq:eic-residual-reflection}
\end{equation}
For the forward direction, substitute $x=Sa$, $y=b$ into
\eqref{eq:eic-residual-invariant}.  If $d(Sa,b)=1$, the sign is $1$
and the selector returns its second branch, giving
\[
 b+1=Sa.
\]
The defining equation for addition gives $b+1=S(b+0)=Sb$.  Thus
$Sb=Sa$, and applying $P$ gives $b=a$, hence $a=b$.
Conversely, $a=b$ and \eqref{eq:eic-diagonal} give
$d(Sa,b)=d(Sa,a)=1$.

Define the equality characteristic by the fixed description
\begin{equation}
 e(a,b)=\operatorname{sg}\bigl(K(d(Sa,b))\bigr).
 \label{eq:eic-equality-characteristic}
\end{equation}
The sign, countdown, and residual reflection lemmas give
\begin{align}
 e(a,b)=0&\quad\longleftrightarrow\quad a=b,
       \label{eq:eic-equality-reflection}\\
 e(a,b)=0&\ \lor\ e(a,b)=1.
       \label{eq:eic-equality-range}
\end{align}
In detail, $e(a,b)=0$ is equivalent to $K(d(Sa,b))=0$ by
\eqref{eq:eic-sign-zero}, which is equivalent to $d(Sa,b)=1$ by
\eqref{eq:eic-countdown-zero}, and then to $a=b$ by
\eqref{eq:eic-residual-reflection}.  The range assertion is
\eqref{eq:eic-sign-range} applied to $K(d(Sa,b))$.

The reflection property concerns the logic's actual equality throughout
the domain.  It was not assumed as part of a numerical comparison symbol,
and its proof did not restrict $a$ or $b$ to standard numerals.
It also yields decidable equality.  In the $e(a,b)=0$ branch,
\eqref{eq:eic-equality-reflection} gives $a=b$.  In the $e(a,b)=1$
branch, an assumption $a=b$ would give $e(a,b)=0$ and therefore $0=1$,
contradicting separation.  Hence
\begin{equation}
 \forall a\,\forall b\,\bigl(a=b\lor(a=b\to\bot)\bigr)
 \label{eq:eic-equality-decidable}
\end{equation}
is a theorem, not an extra hypothesis.

\subsection{Compilation of all quantifier-free formulas}

For each fixed quantifier-free formula $A(\vec p)$, construct a term
$b_A(\vec p)$ by external recursion on its finite syntax.  Take equality,
$\bot$, $\land$, $\lor$, and $\to$ as the primitive formula constructors;
negation is implication into $\bot$.  Set
\begin{align}
 b_{s=t}&=e(s,t), & b_{\bot}&=1,\label{eq:eic-boolean-atoms}\\
 b_{A\land D}&=C(b_A,b_D,1),\label{eq:eic-boolean-and}\\
 b_{A\lor D}&=C(b_A,0,b_D),\label{eq:eic-boolean-or}\\
 b_{A\to D}&=C(b_A,b_D,0).\label{eq:eic-boolean-implies}
\end{align}
Every term here is a finite composition of fixed primitive-recursive
descriptions; if desired, projections and composition give it a single
description symbol.  No internal quantification over formulas or programs
is involved.

An external induction on the syntax of $A$ proves, within the arithmetic
theory,
\begin{equation}
 \bigl(b_A=0\lor b_A=1\bigr)
 \quad\text{and}\quad
 \bigl(b_A=0\longleftrightarrow A\bigr).
 \label{eq:eic-boolean-correctness}
\end{equation}
Here are the full connective checks.  In every case $b_A=1$ implies
$\neg A$: if $A$ held, the already established implication
$A\to b_A=0$ would contradict separation.

\paragraph{Atomic formulas and falsity.}
For an equation, both claims are
\eqref{eq:eic-equality-reflection}--\eqref{eq:eic-equality-range}.
For falsity the range assertion is $1=0\lor1=1$, with its right disjunct
proved by reflexivity.  Also $1=0\to\bot$ follows from separation and
symmetry, while $\bot\to1=0$ follows by explosion.  Thus $1=0$ is
equivalent to $\bot$, as required.

\paragraph{Conjunction.}
Split $b_A=0\lor b_A=1$.  In the zero branch, $A$ holds and
$b_{A\land D}=b_D$.  Its range follows from the range of $b_D$,
and its being zero is equivalent to $D$, hence to $A\land D$ under
the known $A$.  In the one branch, $\neg A$ holds and
$b_{A\land D}=1$.  Its range is immediate.  If it were zero,
separation would give $\bot$, hence $A\land D$; conversely an
assumption $A\land D$ contradicts $\neg A$, and explosion gives
$b_{A\land D}=0$.

\paragraph{Disjunction.}
In the branch $b_A=0$, the formula $A$ holds and
$b_{A\lor D}=0$.  The range is immediate, and both $b_{A\lor D}=0$
and $A\lor D$ hold, proving their equivalence.  In the branch
$b_A=1$, we have $\neg A$ and $b_{A\lor D}=b_D$.  Its range follows
from the range of $b_D$.  If its value is zero, then $D$ and hence
$A\lor D$.  Conversely, from $A\lor D$, the $A$ alternative
contradicts $\neg A$ and yields $b_D=0$ by explosion, while the
$D$ alternative yields $b_D=0$ directly.  This proves the required
equivalence in both alternatives.

\paragraph{Implication.}
In the branch $b_A=0$, the formula $A$ holds and
$b_{A\to D}=b_D$.  Its range follows from the range of $b_D$.
If its value is zero, then $D$, which gives $A\to D$.  Conversely,
$A\to D$ and the known $A$ give $D$, hence $b_D=0$.
In the branch $b_A=1$, we have $\neg A$ and
$b_{A\to D}=0$.  Its range is immediate.  To prove $A\to D$,
assume $A$; this contradicts $\neg A$, and explosion gives $D$.
Thus both $b_{A\to D}=0$ and $A\to D$ hold in this branch, and
their equivalence follows.  This argument uses the proved two-value
range of the finite characteristic.  It does not use excluded middle
for arbitrary formulas or change the surrounding intuitionistic logic.

This completes the external syntactic induction proving
\eqref{eq:eic-boolean-correctness} for every quantifier-free formula.

\subsection{Recovery of induction and the full theory}

Let $A(n,\vec p)$ be any fixed quantifier-free formula.  Suppose an
arbitrary retained context $\Gamma$ proves its base and its induction
step:
\[
 \Gamma\vdash A(0,\vec p),
 \qquad
 \Gamma,A(z,\vec p)\vdash A(Sz,\vec p),
\]
with the usual freshness condition on $z$.  By
\eqref{eq:eic-boolean-correctness}, the base gives
$\Gamma\vdash b_A(0,\vec p)=0$.  Under the additional assumption
$b_A(z,\vec p)=0$, the same equivalence gives $A(z,\vec p)$,
the given step gives $A(Sz,\vec p)$, and the equivalence gives
$b_A(Sz,\vec p)=0$.  Apply literal-equation induction to
$b_A(n,\vec p)=0$ and then apply its reflection implication.
The result is $\Gamma\vdash A(t,\vec p)$ for every target term $t$.
Thus the theory proves the full contextual quantifier-free induction
rule, including in arbitrary quantified contexts.

\paragraph{Completion theorem.}
The same proof first applies to the scalar language, since all the
terms used above have the scalar definitions given at the start of this
appendix. It yields the QFI needed by the subsequent PR compiler.
After that compiler, in the common language of fixed primitive-recursive
descriptions, the theory
\[
 \mathsf{PRE}+\mathsf{Cov}+\mathsf{Sep}+\mathsf{EI}
\]
has exactly the same sequents as ordinary first-order intuitionistic
primitive recursive arithmetic with all quantifier-free induction
instances.  The forward simulation follows from the recovered
quantifier-free induction, successor injectivity, and successor separation;
all primitive-recursive defining equations and all logical rules are
shared.  For the reverse simulation, literal-equation induction is an
instance of quantifier-free induction, and successor separation gives
$\mathsf{Sep}$.  Quantifier-free induction proves $\mathsf{Cov}$:
at zero choose its left disjunct, and at a successor use
$Sx=S(P(Sx))$, which follows from $P(Sx)=x$, for its right disjunct.

Both interpretations therefore use the full domain, defined by $x=x$,
and preserve every function symbol, equality, logical connective, and
quantifier.  Both composites have the identity isomorphism, with graph
$x=y$.  The theories are consequently identity-domain bi-interpretable
and equiconsistent.  This completion theorem is unconditional once the
literal-equation induction rule has been obtained; it has no dependency
on a classical successor-uniqueness reduction.

\section{Pairing, primitive recursion, and internal numerical composites}
\label{sec:iterator-coding}
This appendix is adapted with attribution from the progress-closure
arithmetic development~\cite{progress}. Its prerequisites have now been
proved: scalar predecessor, arithmetic reflection, and full quantifier-free
induction. Pairing is a derived proof tool; it is not part of any checker.

\subsection{Pairing from a bounded square-root search}

Once QFI is available, a familiar square-based injection gives a
particularly direct compiler. Write $a\le b$ for $d(a,b)=0$.
The usual linear-order and semiring identities follow by fixed
quantifier-free inductions from the successor, addition, and monus
equations. Multiplication itself is the scalar term
\[
                        a\cdot b=I(b,0;u.u+a).
\]
Only primitive-recursive arithmetic is involved in the following
bounded root search; there is no minimization over an unbounded domain.

For a fixed parameter $n$, put
\[
 U_n(v)=C(d((Sv)^2,n),Sv,v),\qquad
 R_n(i)=I(i,0;v.U_n(v)),\qquad q(n)=R_n(n).
\]
The update increments $v$ exactly when $(v+1)^2\le n$, and
otherwise retains it. QFI in $i$, with $n$ fixed, proves
\[
 R_n(i)\le i,\qquad R_n(i)^2\le n,\qquad
 R_n(i)=i\ \vee\ n<(S(R_n(i)))^2.              \tag{Root invariant}
\]
The base has $R_n(0)=0$. If the update increments, its test
contradicts the strict-upper-bound alternative in the third invariant.
Thus the old state equals $i$, and the new state equals $Si$;
its square bound is the increment test. If the update does not increment, the
strict upper square bound holds and the state remains fixed.
An already established strict upper square bound forces every
subsequent update to remain fixed. These are exactly the cases
needed for the three displayed invariants.

At $i=n$, either the strict upper square bound already holds or
$R_n(n)=n$, in which case $n<(n+1)^2$ supplies it. Thus
\[
                        q(n)^2\le n<(q(n)+1)^2.           \tag{Root}
\]
Monotonicity of squares shows uniqueness of a number satisfying
Root. All formulas just used are quantifier free in the scalar
language, with fixed $n$ during the induction.

Now define
\[
 \pi(a,b)=(a+b)^2+a,\qquad
 L(n)=d(n,q(n)^2),\qquad H(n)=d(q(n),L(n)).
\]
For $k=a+b$, one has $k^2\le k^2+a<(k+1)^2$ since $a\le k$.
Root uniqueness gives $q(\pi(a,b))=a+b$, and therefore
\[
                    L(\pi(a,b))=a,\qquad H(\pi(a,b))=b.
\]
Surjectivity of $\pi$ is unnecessary. The original numerical domain
is never replaced by the range of this coding; pairing is used only
to store the finite state of a fixed primitive recursion. This
standard square pairing is a convenient reconstruction here, not
a claimed new arithmetic basis.

\subsection{Compiling each primitive recursion}

Suppose the fixed target description is
\[
 f(0,\vec p)=g(\vec p),\qquad
 f(Sn,\vec p)=h(n,f(n,\vec p),\vec p).
\]
Assume by external structural induction that the earlier $g,h$
already have scalar translations $g^*,h^*$. Define
\[
 W(u,\vec p)=\pi(SL(u),h^*(L(u),H(u),\vec p)),
\]
\[
 V(n,\vec p)=I(n,\pi(0,g^*(\vec p));u.W(u,\vec p)),
 \qquad f^*(n,\vec p)=H(V(n,\vec p)).
\]
The update has only its state variable bound, with every $p_i$
fixed. Equation induction proves $L(V(n,\vec p))=n$.
The two projection equations and scalar computation then prove
both defining equations for $f^*$ internally. Initial functions
and composition translate directly. This completes the translation
of every fixed finite PR description without adding a universal
evaluator.

\subsection{Why both numerical round trips are internal}

Two fixed scalar iterators with the same seed and provably equal
updates are internally equal at all counters: use literal-equation
induction comparing their values. Their base agrees, and their
successor equations plus congruence prove the step. The same
uniqueness argument works for ordinary primitive recursion.

For an ordinary PR description, induct externally on its finite
construction. The double translation satisfies its original
defining equations, so internal equation induction identifies it
with the original operation. For a native scalar description,
earlier translated updates agree by the external induction
hypothesis; scalar uniqueness identifies the double-translated
iterator. These proofs hold at arbitrary internal counters.

Thus the domain relation $x=y$ is provably a total bijection
preserving every fixed numerical operation in both composites.
The formula constructions and round trips in
Section~\ref{sec:choice-introduction} give the corresponding formula
composites. Both interpretations preserve falsity, taking native
$0=1$ to an ordinary contradiction and ordinary $\bot$ to native
$0=1$. Finite proof translations therefore establish equiconsistency.


\begin{thebibliography}{9}
\bibitem{shared}VladimirReshetnikov/ProveIt,
\emph{A Shared Continuation for iPRA},
\file{ipra_shared_continuation_2026-11-10.zip}, report section
``A genuinely smaller numerical grammar: the precise unresolved step.''
Pinned corpus commit \file{2fdcf5bf988bead821c416649f713f98b7ed8776}.
\href{https://github.com/VladimirReshetnikov/ProveIt/blob/2fdcf5bf988bead821c416649f713f98b7ed8776/docs/drafts/ipra_shared_continuation_2026-11-10.zip}{Pinned source archive}.
\bibitem{generated}VladimirReshetnikov/ProveIt,
\file{ipra_generated_assumptions_2026-11-10.zip}, code, protocol,
zero-instance proof, and attributed scalar appendices, at the same commit.
\href{https://github.com/VladimirReshetnikov/ProveIt/blob/2fdcf5bf988bead821c416649f713f98b7ed8776/docs/drafts/ipra_generated_assumptions_2026-11-10.zip}{Pinned source archive}.
\bibitem{successor}VladimirReshetnikov/ProveIt,
\file{ipra_successor_recursion_2026-11-09.zip},
\file{theory/successor_recursion.tex}, at the same commit.
\href{https://github.com/VladimirReshetnikov/ProveIt/blob/2fdcf5bf988bead821c416649f713f98b7ed8776/docs/drafts/ipra_successor_recursion_2026-11-09.zip}{Pinned source archive}.
\bibitem{elimination}VladimirReshetnikov/ProveIt,
\file{ipra_elimination_assumptions_2026-11-09.zip}, native logical
macros and formula return translations, at the same commit.
\bibitem{forward}ProveIt forward-left arithmetic development,
as reproduced with attribution in the generated-assumptions archive,
\file{reference/equality_reflection.tex}. Its proof is included here with
the instantiated predecessor reference adjusted.
\bibitem{progress}ProveIt progress-closure arithmetic development,
as reproduced with attribution in the generated-assumptions archive,
\file{reference/primitive_recursion.tex}. Its scalar compiler is included
here with notation local to the appendix.
\bibitem{salehi}S. Salehi, \emph{Decidable formulas of intuitionistic
primitive recursive arithmetic}, Reports on Mathematical Logic
\textbf{36} (2002), 55--61.
\url{https://rml.tcs.uj.edu.pl/rml-36/36-SALEHI.pdf}.
\bibitem{severin}D. E. Severin, \emph{Unary primitive recursive functions},
Journal of Symbolic Logic \textbf{73} (2008), 1122--1138.
See in particular Remarks 5.10 and 6.9 for related seed shifts.
\url{https://arxiv.org/abs/cs/0603063}.
\end{thebibliography}
\end{document}
